\documentclass[conference]{IEEEtran}
\IEEEoverridecommandlockouts
\usepackage{cite}
\usepackage{amsmath,amssymb,amsfonts}
\usepackage{algorithmic}
\usepackage{graphicx}
\usepackage{textcomp}
\usepackage{xcolor}
\usepackage{booktabs}
\usepackage{multirow}
\usepackage{array}
\usepackage{placeins}
\usepackage{epsfig}
\usepackage{tikz}
\usetikzlibrary{positioning,arrows.meta,calc}

\def\BibTeX{{\rm B\kern-.05em{\sc i\kern-.025em b}\kern-.08em
    T\kern-.1667em\lower.7ex\hbox{E}\kern-.125emX}}

\begin{document}

\title{Edge AI-Based FPGA Accelerator for Spectrum Classification and Anomaly Detection}

\author{Author 1, Author 2, and Author 3
  %  \IEEEauthorblockN{Yash~Agrawal, Srijan~Mukhopadhyay, and Sandeep Joshi}
  %  \IEEEauthorblockA{Department of Electrical and Electronics Engineering\\ Birla Institute of Technology and Science, Pilani, Rajasthan 333031, India}
  %  \{f20230306, f20230334, sandeep.joshi\}@pilani.bits-pilani.ac.in
  }
\maketitle


\begin{abstract}
Radio frequency (RF) spectrum monitoring is critical for modern electromagnetic signal environments, where non-cooperative signals, including jamming pulses, spoofing transmissions, and wide-band interference, must be classified and countered within a fraction of a millisecond of first reception. This work presents a fully synthesizable field-programmable gate array (FPGA) accelerator for real-time RF anomaly detection at the sensor edge, eliminating dependence on cloud connectivity or external memory. The system implements an eight-layer one-dimensional residual convolutional neural network (ResNet-1D).
%The trained model achieves 83.26\% test accuracy in 32-bit floating-point (FP32) precision and 83.18\% after 8-bit integer (INT8) symmetric post-training quantization, representing a degradation of only 0.08\% points.
Hardware deployment is enabled through batch normalization (BN) folding and symmetric per-channel INT8 quantization, compressing all weights to 28.6~KB within the block random access memory (BRAM) budget of a Xilinx Zynq-7000 system-on-chip (SoC). The intellectual property (IP) core, governed by a descriptor-driven layer sequencer and closing timing at 142.86~MHz, achieves a compute latency of 142.2~$\mu$s (161.9~$\mu$s measured end to end on the ZedBoard, 60.5$\times$ faster than the on-chip ARM processor) at an estimated accelerator power of only 0.290~W, i.e., 41.2~$\mu$J per inference, which makes it suitable for battery-powered, real-time spectrum monitoring.
\end{abstract}

\begin{IEEEkeywords}
Edge AI, FPGA accelerator, RF anomaly detection, ResNet-1D, Spectrum monitoring.
\end{IEEEkeywords}

\section{Introduction}

Modern electromagnetic spectrum monitoring demands near real-time detection of non-cooperative emissions, such as jamming, spoofing, and wideband interference, within a fraction of a millisecond of reception. FPGAs provide an ideal edge platform by synthesizing dedicated convolutional neural network (CNN) datapaths that avoid the instruction overheads and cache misses of central processing units (CPUs) and graphics processing units (GPUs), delivering high spatial parallelism and deterministic sub-millisecond latency \cite{b1, b2}. However, edge FPGA deployment is heavily constrained by the steep look-up table (LUT) and digital signal processing (DSP) costs of floating-point operations and runtime batch normalization (BN) statistics, necessitating quantization and structural folding.
%The problem of all these three limitations is overcome in an overall hardware-software co-design approach.

In recent years, end-to-end learning on raw in-phase (I) and quadrature (Q) samples has become the mainstream approach in automatic modulation classification (AMC) research \cite{b1, b2}. %The last decade surveys have deeply summarized the AMC systems development process from the earliest applications of CNNs to the recent designs of transformer and hybrid networks~\cite{b1}. A review of recent works~\cite{b2} concerning the deeper insertion of DL for physical-layer communication tasks like channel estimation, signal detection, and modulation classification, has motivated hardware-accelerated edge inference.
The RadioML 2018.01A dataset \cite{b3} provides 24-class I/Q data at 26 signal-to-noise ratio (SNR) levels and has become the community-standard benchmark. The authors in \cite{b4} demonstrated the efficiency of an attention-enhanced, transformer-based deep learning architecture. The authors in \cite{b5} have developed a systematic description of robust deep learning for physical layer communications.
%Hardware acceleration of RF classification has followed two main streams: CPU/GPU-based offline processing, and FPGA-based edge inference. GPU inference of a full RadioML classifier typically requires 5--20~ms per batch on an NVIDIA Jetson-class device, impractical for sub-microsecond EW reaction times.
Recent literature on FPGA-based AI accelerators \cite{b6} categorizes optimizations at the computational, memory, and system levels, emphasizing custom register transfer level (RTL) and extreme quantization as key for performance in resource-limited systems. The authors in \cite{b7} showed that attention-based models achieve better performance on wireless modulation classification.
Existing FPGA deep learning RF classifiers, such as the DeepRadio platform \cite{b8}, classify signals received through the RF front end into different modulation types in real time with low power. The RF System-on-Chip (RFSoC) work in \cite{b9} uses VGG10 and ResNet33 models on the RadioML 2018.01A dataset with FPGA deployment. While prior work in \cite{b10} demonstrated an FPGA-accelerated anomaly detector utilizing an unsupervised autoencoder for generic, out-of-distribution flagging, our work distinctly frames the problem around a supervised ResNet-1D architecture mapping directly to a robust, operational signal taxonomy. The authors in~\cite{b11} achieved a high classification accuracy using a ResNet with six residual stacks on the full 24-class RadioML dataset. However, the model in~\cite{b11} exceeds edge-FPGA BRAM capacity, requiring high-latency off-chip memory, while fine-grained modulation labels lack operational utility and cannot distinguish legitimate traffic, jamming, spoofing, and co-channel interference.
%Tools like hls4ml already implement optimizers that fuse batch normalizations following affine projection layers into a single node with equivalent fused weights \cite{b12}, and ResNet INT8 inference on Xilinx devices is established practice. Recently, \cite{b13} surveyed deep transfer learning approaches for AMC, highlighting emerging architectures that improve generalization across diverse channel conditions. The authors in \cite{b14} demonstrated an FPGA-based real-time signal measurement system using deep learning accelerators.
%However, this work contributes a genuinely novel element that is absent in prior literature: a defense-threat taxonomy that remaps the 24 RadioML classes into four EW-operational categories (Normal, Jammer, Spoofer, Interference), combined with a fully hand-written Verilog RTL implementation with zero High-Level Synthesis (HLS) tooling dependencies. No prior work performs this exact threat-aligned remapping with a complete custom RTL datapath.

%The limitations motivate an end-to-end hardware-accelerated pipeline implemented on an FPGA in this work; the major contributions are summarized as follows.

Motivated by the limitations of fine-grained classifiers, excessive parameter footprint, and edge detectors lacking threat classification, in this work, we propose a fully synthesizable FPGA accelerator for real-time operational signal classification and anomaly detection, through a hardware-software co-design: we frame the problem as an actionable operational threat taxonomy, compress an 8-layer ResNet-1D to under 30~KB using integer-only quantization with scale-aligned residual additions, and implement a parameterizable processing element (PE) datapath in hand-written Verilog register transfer level (RTL) code that executes entirely from on-chip BRAM, with no external memory accesses, on a Xilinx Zynq-7000 FPGA. The main contributions of this work are as follows:
\begin{itemize}
    \item We introduce a novel four-class operational signal taxonomy that maps the 24 RadioML 2018.01A modulation classes onto operational signal categories (Normal, Jammer, Spoofer, Interference), reframing fine-grained modulation recognition as an operationally meaningful signal-classification problem.
    \item We design an eight-layer ResNet-1D architecture, optimized for FPGA deployment through symmetric post-training INT8 quantization and algebraic BN folding, compressing all weights to 28.6~KB of on-chip BRAM. The integer-only INT8 model preserves 83.18\% test accuracy, only 0.08 percentage points below the FP32 baseline.
    \item We design and implement a ResNet-1D accelerator IP core (RadarNet-1D) that runs bit-exactly with the integer model, achieving a measured end-to-end latency of 161.9~$\mu$s (60.5$\times$ faster than the on-chip ARM Cortex-A9) and an estimated accelerator power of 0.290~W, while utilizing just 5.86\% of LUTs, 2.30\% of flip-flops, 10.00\% of BRAM, and 47.27\% of DSP slices on the Xilinx Zynq-7000 ZedBoard.
    %We functionally verify the RTL against 50 Python-generated golden test vectors with per-vector PASS/FAIL reporting and cycle-accurate latency benchmarking, and validate the design on a Xilinx Zynq-7000 ZedBoard using ILA/VIO-based hardware debug, confirming
\end{itemize}
%The remainder of this paper is organized as follows. Section~II presents the proposed architecture; Section~III gives the details of the FPGA implementation, Section~IV gives the results, and Section~V concludes the work.
\section{Proposed Architecture}
The RadioML 2018.01A dataset \cite{b3} provides 2,555,904 labeled I/Q signal records across 24 modulation types at 26 SNR levels from $-20$~dB to $+30$~dB in 2~dB steps. Each record is a floating-point array of shape $(1024, 2)$ representing 1,024 complex-valued time samples with I/Q baseband components. The dataset models realistic channel impairments including additive white Gaussian noise (AWGN), multipath Rayleigh fading, carrier frequency offset, timing recovery errors, and hardware non-linearities such as I/Q imbalance and phase noise. In this work, only records with SNR $\geq 0$~dB are retained, the first 128 samples of each record are used as the input frame, and 50,000 frames per class are drawn and split 70/15/15 into training, validation, and test sets (30,000 test frames).
Standard modulation recognition assigns one label per modulation type, yielding a 24-class problem with fine-grained distinctions, e.g., between 16-ary and 32-ary quadrature amplitude modulation (16-QAM and 32-QAM). For real-time spectrum monitoring applications, this granularity is operationally unnecessary: it increases classifier complexity while the system only needs to know whether a signal is benign or what operational signal category it represents. This work instead introduces a four-class operational signal taxonomy that consolidates the 24 modulations into operationally aligned groups based on spectral characteristics, bandwidth occupancy, and operational signal profiles.
% TODO(authors): JUSTIFY THE TAXONOMY BEFORE SUBMISSION.
% This is the paper's headline novelty claim, but Table I currently presents the Normal/Jammer/Spoofer/
% Interference mapping with no citation and no domain justification beyond the one-line "Rationale"
% column. A reviewer with EW/SIGINT background is likely to challenge specific assignments -- e.g. OOK
% is a simple, commonly legitimate modulation, so labeling it "Jammer" needs defending. Recommend either
% (a) citing EW/SIGINT literature or doctrine that supports grouping by these spectral/bandwidth
% properties, or (b) explicitly framing the taxonomy as an illustrative/candidate mapping for
% demonstrating hardware feasibility, with real-world taxonomy validation noted as future work.
% \begin{table}[htbp]
% \caption{Defense-Threat Class Mapping from RadioML 2018.01A Modulations}
% \begin{center}
% \footnotesize
% \begin{tabular}{c l p{2.2cm} p{2.5cm}}
% \toprule
% \textbf{ID} & \textbf{Category} & \textbf{Rationale} & \textbf{RadioML Modulations} \\
% \midrule
% 0 & Normal & Narrowband continuous-wave, legitimate emissions & FM, GMSK, AM-DSB-WC, AM-DSB-SC, AM-SSB-WC, AM-SSB-SC \\
% \midrule
% 1 & Jammer & Wideband high-entropy, consistent with barrage jamming & OOK, 64QAM, 128QAM, 256QAM \\
% \midrule
% 2 & Spoofer & Coherent phase-modulated, may mimic radar returns & BPSK, QPSK, 8PSK, 16PSK, 32PSK, OQPSK \\
% \midrule
% 3 & Interference & Multi-level amplitude-phase, cross-band interference & 16APSK, 32APSK, 64APSK, 128APSK, 16QAM, 32QAM, 8ASK, 4ASK \\
% \bottomrule
% \end{tabular}
% \label{tab:threat_mapping}
% \end{center}
% \footnotesize{MPSK: M-ary Phase Shift Keying, MQAM: M-ary Quadrature Amplitutde Modeling, FM: Frequency Modulation, GMSK: Gaussian Minimum Shift Keying, AM-DSB-WC: Amplitude Modulation, Double Sideband, With Carrier,AM-DSB-SC: Amplitude Modulation, Double Sideband, Suppressed Carrier,AM-SSB-WC: Amplitude Modulation, Single Sideband, With Carrier, AM-SSB-SC: Amplitude Modulation, Single Sideband, Suppressed Carrier, OOK: On-Off Keying, BPSK: Binary Phase shift keying, QPSK: Quadrature Phase shift Keying, OQPSK: Offset Quadrature Phase Shift Keying, MASK: M-ary Amplitude Shift Keying, MAPSK: M-ary Amplitude Phase Shift Keying}
% \end{table}

A plain stacked 1D CNN applies convolutional layers sequentially with no bypass paths. During backpropagation, the gradient of the loss with respect to early-layer parameters is computed as a product of Jacobian matrices across every intermediate layer, which attenuates the gradient as depth grows. Residual shortcut connections add a block's input directly to its output, so the gradient bypasses these products and reaches early layers intact; the proposed network therefore uses two residual blocks.
%When these matrices contain values consistently less than 1, the product shrinks exponentially with depth according to the vanishing gradient problem \cite{b15,b16}. In practice, plain CNNs beyond 8--10 layers learn little in their early stages and often perform worse than shallower counterparts despite having more parameters. Residual Networks \cite{b14} solve this by adding a shortcut connection that adds the block's input directly to its output, bypassing one or two convolutional layers. Gradients flow directly backward through the shortcut without attenuation, reaching early layers intact regardless of network depth. A further hardware advantage of ResNet-1D is that Batch Normalization layers can be algebraically folded into the preceding convolution weights prior to quantization, eliminating every BN-related computation \cite{b16}.
ResNet-1D maps a two-channel 128-point I/Q input to four class logits through eight weight layers (seven convolutions and one fully connected layer). The architecture employs progressively increasing channel widths ($2 \rightarrow 16 \rightarrow 32 \rightarrow 64$) combined with stride-2 downsampling at two points to compress the temporal dimension from 128 to 32 timesteps before global average pooling. Table~\ref{tab:layerwise} provides a complete layer-wise breakdown.
%%%%%%%%%%%%%%%%%%%%%%%%%%%%%%%
\begin{table}[t]
\renewcommand{\arraystretch}{1.15}
\caption{ResNet-1D Layer-wise Architecture}
\label{tab:layerwise}
\centering
\scriptsize
\setlength{\tabcolsep}{3pt}
\begin{tabular}{c l c c c r}
\toprule
\textbf{Layer} & \textbf{Type} & \textbf{In (C$\times$L)} & \textbf{Out (C$\times$L)} & \textbf{k, s} & \textbf{Parameters} \\
\midrule
1 & Conv + BN + ReLU & $2\!\times\!128$ & $16\!\times\!128$ & 7, 1 & 240 \\
2 & Conv + BN + ReLU & $16\!\times\!128$ & $32\!\times\!64$ & 3, 2 & 1,568 \\
3--4 & ResBlk 1 (id) & $32\!\times\!64$ & $32\!\times\!64$ & 3, 1 & 9,280 \\
5--7 & ResBlk 2 (proj) & $32\!\times\!64$ & $64\!\times\!32$ & 3, 2 & 17,664 \\
--- & GAP & $64\!\times\!32$ & $64\!\times\!1$ & --- & 0 \\
8 & FC + Argmax & 64 & 4 & --- & 260 \\
\midrule
 & \multicolumn{4}{r}{\textbf{Total Trainable Parameters}} & \textbf{29,252} \\
\bottomrule
\end{tabular}

\vspace{2pt}
\parbox{\columnwidth}{\scriptsize Conv: 1-D convolution; ResBlk: residual block; id: identity shortcut; proj: $1\times1$ projection shortcut; GAP: global average pooling; FC: fully connected layer; ReLU: rectified linear unit; C: channels; L: temporal length; k: kernel size; s: stride.}
\end{table}
%%%%%%%%%%%%%%%%%%%%%%%%%%%%%%%
% [FIX] "indice" → "index" (two occurrences below)
%The datapath is governed by the following equations, derived directly from the layer definitions in Table~\ref{tab:layerwise} and implemented one-to-one in the RTL described in Section~\ref{sec:hw}.
The 1D convolution output for output channel index $c_o$ at time position $t$ is expressed as
\begin{equation}
y[c_o, t] = \sum_{c_i=0}^{C_{in}-1} \sum_{k=0}^{K-1} W[c_o, c_i, k]\, x[c_i, t s + k - p] + b[c_o],
\label{eq:conv1d}
\end{equation}
where $c_i$ is the input channel index, $K$ is the kernel size, $s$ is the stride, $p$ is the zero-padding width, $x$ is the input tensor, $W$ is the weight tensor, $[ \cdot, \cdot,\cdot]$ denotes tensor indexing, $C_{in}$ is the number of input channels, and $b[c_o]$ is the bias. BN normalizes each channel, and the normalized output for each channel is expressed as
\begin{equation}
\hat{x}[c] = \gamma[c] \frac{x[c] - \mu[c]}{\sigma[c]} + \beta[c],
\label{eq:bn}
\end{equation}
where $x[c]$ is the input for channel $c$, $\mu[c]$ and $\sigma[c]$ are the channel's running mean and standard deviation, and $\gamma[c]$ and $\beta[c]$ are the learned scale and shift parameters, respectively. ReLU activation applies element-wise, $\text{ReLU}(x) = \max(0, x)$, where $x$ is the input value and $\max$ returns the maximum of its arguments. The identity skip connection computes $y = F(x) + x$ when channel dimensions match; when they do not, a $1\times1$ projection convolution adjusts the identity as $y = F(x) + W_p\ast x$, where $x$ is the input tensor, $y$ is the output tensor, $F(x)$ is the residual mapping function, and $W_p$ is the $1\times1$ projection weight matrix. Global Average Pooling reduces each channel to a scalar. The pooled output scalar for channel $c$ is given as $ z[c] = \frac{1}{T} \sum_{t=0}^{T-1} x[c, t],$
where $T$ is the temporal sequence length, $t$ is the time step index, and $x[c, t]$ is the input value at channel $c$ and time step $t$. Since $T = 32$, the hardware division is replaced by a rounded arithmetic right shift by 5 bits, $(\sum_t x[c,t] + 16) \gg 5$, eliminating all division circuitry from the pooling datapath. The argmax classification selects $\hat{c} = \arg\max_{j \in \{0,1,2,3\}} l[j]$, and an anomaly flag is asserted whenever $\hat{c} \neq 0$, where $\hat{c}$ is the predicted class index, $\arg\max$ is the operation returning the argument that gives the maximum value, $j$ is the class index iterating over the four classes, and $l[j]$ is the logit score for class $j$.

\section{INT8 Quantization and FPGA Implementation}
%%%%%%%%%%%%%%%%%%%%%%%%%%%%%%%%
\begin{figure*}[t]
\centering
\begin{minipage}[t]{0.49\textwidth}
  \centering
  \includegraphics[width=\textwidth]{radarnet_hw_overview.pdf}
  \caption{Top-level datapath of the proposed RadarNet-1D accelerator.}
  \label{fig:hw_overview}
\end{minipage}\hfill
\begin{minipage}[t]{0.49\textwidth}
  \centering
  \includegraphics[width=\textwidth]{radarnet_dsp_mac_engine.pdf}
  \caption{One of the $W$ identical output lanes of the PE array.}
  \label{fig:dsp_mac}
\end{minipage}
\end{figure*}
%%%%%%%%%%%%%%%%%%%%%%%%%%
% [FIX] Restored missing opening sentence with correct citation b12 (Jacob, quantization)
\begin{figure}[t]
\centering
\includegraphics[width=\columnwidth, height=2in]{Fig1.jpeg}
\caption{Proposed edge-AI deployment setup showing the ZedBoard running the INT8 ResNet-1D accelerator for anomaly detection.}
\label{fig:impl}
\end{figure}
This work employs an integer-only quantization scheme based on
symmetric per-tensor activation scaling and per-channel weight
scaling~\cite{b12}. Symmetric quantization keeps every multiply-accumulate (MAC) as a single signed multiplication with no offset accumulation. Before quantization, each BN layer is algebraically absorbed into
its previous convolution, and the folded weight for output channel $c_o$, input channel $c_i$, and kernel tap $k$ is given as
\begin{equation}
\tilde{W}[c_o, c_i, k]
  = \frac{\gamma[c_o]}{\sqrt{\sigma^{2}[c_o]+\epsilon}}\;
    W[c_o, c_i, k],
\label{eq:bnfold}
\end{equation}
where $\sigma^{2}[c_o]$ is the running variance, and $\epsilon$ is a numerical stability constant. The folded bias is computed analogously, absorbing the original bias, running mean $\mu[c_o]$, and BN shift $\beta[c_o]$, eliminating all six BN layers from the RTL. Each folded kernel is then quantized to signed 8-bit precision with per-channel scaling:
\begin{equation}
W_q[c_o,c_i,k]
  = \text{clamp}\!\left(
      \left\lfloor
        \frac{\tilde{W}[c_o,c_i,k]}{s_W[c_o]}
      \right\rceil,\;-128,\;127\right),
\label{eq:quant}
\end{equation}
where $\lfloor\cdot\rceil$ denotes round-to-nearest-integer, $\text{clamp}(v,a,b)$ restricts $v$ to $[a,b]$, and the per-channel scale factor is $s_W[c_o]=\max|\tilde{W}[c_o,:,:]|/127$. Activations are quantized per tensor, with scales calibrated at the 99.99th percentile over 4,096 training frames. INT8$\times$INT8 products are accumulated in a 24-bit register, which bounds the worst-case dot product ($C_{in}K = 192$ terms, $|acc| < 192 \cdot 2^{14} < 2^{22}$), and the INT32 bias is added in the requantization stage. Requantization back to INT8 uses precomputed fixed-point multipliers $M[c_o]<2^{15}$ and arithmetic right-shifts, with a dedicated scale-alignment stage that fuses the residual shortcut addition into a single integer pipeline without any floating-point operation. Fig.~\ref{fig:hw_overview} shows the proposed top-level datapath, including the descriptor-driven sequencer, triple feature buffers, weight read-only memory (ROM), and the classifier/anomaly-detection output stage. Fig.~\ref{fig:dsp_mac} shows the microarchitecture of a single PE output lane of the $W\!\times\!W$ INT8 MAC array, highlighting the DSP48 multipliers, registered adder tree, INT24 accumulator, and the six-stage requantization pipeline (bias, residual add, shift-based rounding, ReLU/saturation). Fig.~\ref{fig:impl} presents the physical deployment on the ZedBoard together with its post-implementation Vivado floorplan, confirming that the INT8-quantized design maps entirely to the Zynq-7000 fabric without routing overflow. The co-design flow spans five stages: Python preprocessing of INT8 test vectors, PyTorch BN folding and quantization, weight export into a $462\times512$-bit ROM, cycle-accurate RTL simulation against golden vectors, and Vivado synthesis and implementation of the accelerator IP on the ZedBoard, with a Vitis bare-metal application on the ARM core that drives it over AXI4-Lite. % For the worst-case layer ($C_{\text{in}} K = 32 \times 3 = 96$), the maximum accumulation magnitude is $96 \times 127^2 = 1.55 \times 10^6 < 2^{21}$, guaranteeing that accumulator registers never overflow 32-bit capacity.
%This work employs symmetric per-tensor post-training quantization (PTQ) \cite{b12}. Asymmetric quantization requires a zero-point subtraction in every MAC operation, expanding to four terms per multiply-accumulate and doubling the effective multiply count in hardware. Symmetric quantization ($z = 0$) keeps the MAC as a single signed multiply with no offset accumulation, saving significant LUT and DSP resources. For a floating-point weight tensor $W$, the quantization scale factor $s_W$ maps the maximum absolute weight value to the INT8 ceiling of 127, given as $s_W = \frac{\max(|W|)}{2^{b-1} - 1},$
%where $b$ is the quantization bit-width, $\max$ denotes the maximum function, and $|W|$ denotes the element-wise absolute value of the weight tensor $W$. Each weight element is quantized by rounding to the nearest integer and clamping to $[-128, 127]$. The resulting quantized weight tensor is expressed as
%\begin{equation}
%W_q = \text{clamp}\left(\text{round}\left(\frac{W}%{s_W}\right), -128, 127\right),
%\label{eq:quant}
%\end{equation}
%where $\text{round (.)}$ is the round-to-nearest-integer operation, and $\text{clamp}(v, \min, \max)$ restricts the value $v$ to the range $[\min, \max]$.
%Prior to quantization, every BatchNorm layer is algebraically absorbed into its preceding convolution. The folded weight is computed per output channel $c_o$ as
%\begin{equation}
%\tilde{W}[c_o, c_i, k] = \frac{\gamma[c_o]}{\sigma[c_o]} W[c_o, c_i, k],
%\label{eq:bnfold_w}
%\end{equation}
%where $c_i$ denotes the input channel,  $k$ denotes the kernel index, $W[c_o, c_i, k]$ is the original convolution weight, and $\gamma[c_o]$ and $\sigma[c_o]$ are the BatchNorm scale and standard deviation for output channel $c_o$. The folded bias absorbs both the original convolution bias and the complete BN affine transformation. The folded bias for output channel $c_o$ is expressed as
%\begin{equation}
%\tilde{b}[c_o] = \frac{\gamma[c_o] \cdot (b[c_o] - \mu[c_o])}{\sigma[c_o]} + \beta[c_o],
%\label{eq:bnfold_b}
%\end{equation}
%where $b[c_o]$ is the original convolution bias, and $\mu[c_o]$ and $\beta[c_o]$ are the BatchNorm running mean and shift parameter for output channel $c_o$. Applying this across all six Conv--BN pairs in ResNet-1D removes every running-mean accumulator, running-variance accumulator, and per-channel divider from the hardware, reducing RTL complexity significantly.

%Each MAC operation multiplies a signed INT8 weight by a signed INT8 input activation, producing a signed 16-bit product. Products are accumulated across all $C_{in} K$ terms into a signed 24-bit register. The maximum accumulation magnitude for the worst-case layer ($C_{in} K = 32 \times 3 = 96$ terms) is $96 \times 127^2 = 1,548,384$, requiring only 21 bits with sign, which is well within the 24-bit accumulator capacity of 8,388,608. After the full dot product and bias addition, the result is saturated back to INT8: $\text{out}_8 = \text{clamp}(\text{acc}_{24} + b_q, -128, 127)$, where $\text{out}_8$ is the saturated 8-bit output, $\text{acc}_{24}$ is the 24-bit accumulated dot product, and $b_q$ is the quantized folded bias.
%This model uses PTQ rather than quantization-aware training (QAT) because QAT incorporates simulated quantization noise during the training loop, allowing the network to adapt its weight distribution to the quantization grid.
%However, QAT requires retraining with a modified loss landscape, adds significant training complexity, and may overfit to the quantization scheme.
%\section{FPGA Implementation}
%\label{sec:hw}
%%%%%%%%%%%%%%%%%%%%%%%%%%%%%%%%
%A top-level module implements the eight-state Moore FSM shown in Fig.~\ref{fig:fsm}, clocked at 142~MHz and targeting Xilinx 7-Series FPGA platforms. In state \textbf{IDLE}, wait for start pulse, which indicates the start of the process, clear busy flag, which indicates if the FSM is doing some operation in a state, de-assert done, which is used to signal the process is over after reaching the final state, res\_valid (result valid) are all set to low. After starting, state \textbf{LOAD} stores 128 I/Q data pairs whenever iq\_valid is high. State \textbf{CONV\_IN} starts convolution using the first two layers, waits until conv\_in\_done, which signals the end of the convolution operation, and becomes true.  In states \textbf{RES\_BLK1} and \textbf{RES\_BLK2}, we implement respectively residual blocks with identity skip and $1 \times 1$ projection skip, waiting until res\_blk1\_done and res\_blk2\_done, which indicate the completion of the above-mentioned operations, become true. In state \textbf{GAP}, global average pooling is executed, and gap\_done is asserted when this execution is completed. In state \textbf{FC\_OUT}, $64 \times 4$ matrix multiplication is carried out, then the argument of the maximum function is run, which finds the index of the largest value in a list. Finally, anomaly\_flag is set to high if the prediction is not class 0.
%In state \textbf{DONE}, capture class\_id, raise done and res\_valid signal to report completion, save total inference cycle, then return to \textbf{IDLE}.
%Three key optimizations distinguish this accelerator: (i)~BN folding removes all six BatchNorm layers from the RTL, eliminating running-statistic accumulators and per-channel dividers, (ii)~GAP replaces division by $T\!=\!32$ with a 5-bit arithmetic right shift, consuming zero multiplier cells, and (iii)~INT8 symmetric quantization compresses all weights to 28.6~KB, fitting within a single 36~KB BRAM block.
%Specifically, an Integrated Logic Analyzer (ILA) core triggers on the \texttt{start} pulse to capture cycle-level waveform data on the physical FPGA. %(Fig.~\ref{fig:ila_trigger}), while
%A Virtual Input/Output (VIO) core provides real-time software dashboard control over inputs like \texttt{iq\_data} and probes state outputs such as \texttt{anomaly\_flag} and \texttt{result\_valid}.

% Placeholder for ILA Trigger figure
%\begin{figure}[htbp]
%\centerline{\includegraphics[width=\columnwidth]{ilatrigger_updt.png}}
%\caption{ILA trigger on inference start, showing hardware execution state transition matching behavioral simulation.}
%\label{fig:ila_trigger}
%\end{figure}

%\begin{figure*}[t]
%\centering
%\includegraphics[width=7in, height=1in]{ilatrigger_updt.png}
%\caption{ILA trigger on inference start, showing hardware execution state transition matching behavioral simulation.}
%\label{fig:ila_trigger}
%\end{figure*}

% Placeholder for VIO Dashboard figure
%\begin{figure}[htbp]
%\centerline{\includegraphics[width=\columnwidth]{viodashboard.png}}
%\caption{VIO Dashboard probing hardware registers. The anomaly flag correctly asserts when anomalous I/Q patterns are injected.}
%\label{fig:vio_dashboard}
%\end{figure}
\begin{table}[t]
\renewcommand{\arraystretch}{1.15}
\caption{Per-Class Performance on Test Set}
\label{tab:perclass}
\centering
\footnotesize
\begin{tabular}{lcccc}
\toprule
\textbf{Class} & \textbf{Precision (\%)} & \textbf{Recall (\%)} & \textbf{F1 (\%)} & \textbf{Support} \\
\midrule
Normal      & 99.22 & 99.45 & 99.33 & 7,500 \\
Jammer      & 70.45 & 68.11 & 69.26 & 7,500 \\
Spoofer     & 93.56 & 96.87 & 95.19 & 7,500 \\
Interference& 68.93 & 68.61 & 68.77 & 7,500 \\
\midrule
\textbf{Overall} & \textbf{83.04} & \textbf{83.26} & \textbf{83.14} & \textbf{30,000} \\
\bottomrule
\end{tabular}
\end{table}
%%%%%%%%%%%%%%%%%%%%%%%%%%
%\subsection{Quantization Pipeline Impact}
%%%%%%%%%%%%%%%%%%%%%%%%%%%%%
\begin{table}[t]
\renewcommand{\arraystretch}{1.15}
\caption{Test Accuracy at Each Stage of the Quantization Pipeline}
\label{tab:quant}
\centering
\footnotesize
\begin{tabular}{lcc}
\toprule
\textbf{Pipeline Stage} & \textbf{Accuracy (\%)} & \textbf{$\Delta$ (\%)} \\
\midrule
FP32 Baseline  & 83.26 & --- \\
After BN Folding (FP32) & 83.26 & 0.00 \\
After INT8 Quantization (integer-only) & 83.18 & $-$0.08 \\
\bottomrule
\end{tabular}
\end{table}

\section{Results}
This section evaluates the proposed system along two axes: classification performance of the ResNet-1D model under the four-class operational signal taxonomy, and the resource, timing, and power characteristics of the synthesized RTL accelerator on the Xilinx Zynq-7000 ZedBoard platform.
%\subsection{Training and Classification Performance}
The ResNet-1D model was trained using PyTorch with the Adam optimizer \cite{b13} at an initial learning rate of $1 \times 10^{-3}$ and weight decay $1 \times 10^{-4}$. A cosine annealing learning rate schedule reduced the rate to $\eta_{\min} = 1 \times 10^{-6}$ over a 60-epoch budget. Early stopping with patience 10 terminated training at epoch 45. Cross-entropy loss was used with a mini-batch size of 256.

Table~\ref{tab:perclass} reports per-class precision, recall, and F1-score on the held-out test set of 30,000 balanced samples. Normal signal classification achieves near-perfect F1 (99.33\%) due to the distinct narrowband spectral characteristics of continuous-wave modulations. Spoofer classification achieves 95.19\% F1, benefiting from the coherent phase structure of PSK modulations. Jammer (69.26\% F1) and Interference (68.77\% F1) exhibit lower scores, attributable to spectral overlap between wideband modulations assigned to different signal classes. Table~\ref{tab:quant} tracks test accuracy at each stage of the quantization pipeline. BN folding (\ref{eq:bnfold}) is algebraically exact and introduces no accuracy change within floating-point precision. Integer-only INT8 quantization introduces a 0.08 percentage-point degradation, well within the design target of less than 1\% and consistent with the rounding bound of $\pm 0.5 s_W$ per weight element \cite{b12}. Table~\ref{tab:power} gives Vivado's estimated on-chip power for the complete ZedBoard system at 142.86~MHz: 1.978~W in total, of which the ARM processing system (PS7) accounts for 1.530~W and the RadarNet-1D accelerator in the programmable logic (PL) for 0.290~W.
%%%%%%%%%%%%%%%%%%%%%%%%%%
%%%%%%%%%%%%%%%%%%%%%%%%%%%%%%%%%%%
%\subsection{Vivado Synthesis and Power Results}
%Vivado synthesis targeting the Xilinx Zynq-7000 SoC available on the ZedBoard platform results in the power numbers presented in Table~\ref{tab:power} and resource utilization in Table~\ref{tab:utilization}.
%%%%%%%%%%%%%%%%%%%%%%%%%%%%%%%
\begin{figure}[t]
\centering
\includegraphics[width=\columnwidth, height=1.8in]{fig4.png}
\caption{Behavioral simulation showing incrementing cycle\_count every inference, class\_id, I/Q data streaming, and control signals.}
\label{fig:waveform}
\end{figure}
% Placeholder for Power Analysis
\begin{table}[t]
\renewcommand{\arraystretch}{1.15}
\caption{On-Chip Power and Latency Breakdown on Xilinx Zynq-7000 ($W{=}8$, 142.86~MHz)}
\label{tab:power}
\centering
\footnotesize
\begin{tabular}{lrr}
\toprule
\textbf{Parameter} & \textbf{Value} & \textbf{Share (\%)} \\
\midrule
\textbf{Total On-Chip Power} & \textbf{1.978 W} & --- \\
\quad Dynamic Power & 1.827 W & 92 \\
\quad Device Static Power & 0.151 W & 8 \\
\midrule
\multicolumn{3}{l}{\textit{Dynamic Power Breakdown}} \\
\quad Processing System (PS7) & 1.530 W & 84 \\
\quad DSP & 0.120 W & 7 \\
\quad Signals & 0.062 W & 3 \\
\quad BRAM & 0.059 W & 3 \\
\quad Logic & 0.034 W & 2 \\
\quad Clocks & 0.022 W & 1 \\
\midrule
\quad \textit{of which RadarNet-1D accelerator (PL)} & \textit{0.290 W} & \textit{16} \\
%\midrule
%Junction Temperature & 47.8$^{\circ}$C & --- \\
%Thermal Margin & 37.2$^{\circ}$C (3.2 W) & --- \\
\midrule[\heavyrulewidth]
\textbf{Latency Component} & \textbf{Cycles} & \textbf{Latency ($\mu$s)} \\
\midrule
I/Q frame load (128 samples) & 128 & 0.90 \\
Layer 1: Conv ($k{=}7$) & 1,817 & 12.72 \\
Layer 2: Conv ($k{=}3$, $s{=}2$) & 1,561 & 10.93 \\
ResBlk 1: Conv 1 & 3,097 & 21.68 \\
ResBlk 1: Conv 2 + identity add & 3,097 & 21.68 \\
ResBlk 2: $1\times1$ projection & 1,049 & 7.34 \\
ResBlk 2: Conv 1 & 3,097 & 21.68 \\
ResBlk 2: Conv 2 + projection add & 6,169 & 43.18 \\
GAP & 263 & 1.84 \\
FC + argmax & 37 & 0.26 \\
\midrule
\textbf{Compute latency (RTL simulation)} & \textbf{20,315} & \textbf{142.2} \\
\quad + AXI-paced I/Q load on the board & +2,698 & +18.9 \\
\quad Accelerator cycle counter (board) & 23,013 & 161.1 \\
\quad + PS-side start and result readout & --- & +0.8 \\
\textbf{End to end, measured on ZedBoard} & --- & \textbf{161.9} \\
\bottomrule
\end{tabular}

\vspace{2pt}
\parbox{\columnwidth}{\scriptsize Power: Vivado vectorless estimate of the complete system. Latency: each layer includes one set-up cycle and a 24-cycle pipeline drain; 1~cycle $=$ 7~ns.}
\end{table}
%%%%%%%%%%%%%%%%%%%%%%%%%%%%%%
%%%%%%%%%%%%%%%%%%%%%%%%%%%%%%
% Placeholder for Resource Utilization
\begin{table}[t]
\renewcommand{\arraystretch}{1.15}
\caption{FPGA Resource Utilization on Xilinx Zynq-7000}
\label{tab:utilization}
\centering
\footnotesize
\begin{tabular}{lccc}
\toprule
\textbf{Resource} & \textbf{Used} & \textbf{Available} & \textbf{Utilization (\%)} \\
\midrule
LUT & 3,117 & 53,200 & 5.86 \\
FF & 2,448 & 106,400 & 2.30 \\
BRAM & 14 & 140 & 10.00 \\
DSP & 104 & 220 & 47.27 \\
%IO & 0 & 200 & 0.00 \\
\bottomrule
\end{tabular}
\end{table}


\begin{table}[t]
\renewcommand{\arraystretch}{1.15}
\caption{Comparison with FPGA RF Classifiers and Anomaly Detectors}
\label{tab:comparison}
\centering
\scriptsize
\setlength{\tabcolsep}{1.2pt}
\begin{tabular}{lccc}
\toprule
\textbf{Metric} & \textbf{Proposed Accelerator} & \textbf{Moss et al. \cite{b10}} & \textbf{FINN~\cite{b14}} \\
\midrule
\textbf{Architecture}   & ResNet-1D (8 layers)    & Autoencoder & Binarized NN (BNN) \\
\textbf{Dataset}        & RadioML 2018.01A        & Custom RF / QPSK      & RadioML \\
\textbf{Quantization}   & INT8 integer-only & Fixed-point 16-bit    & 1-bit (Binary) \\
%\textbf{RTL Design}     & \textbf{Custom Verilog} & Custom RTL & Vivado HLS \\
\textbf{Signal Taxonomy}& \textbf{4-class Operational} & Generic Anomaly  & 10-class Modulation \\
\textbf{Accuracy (\%)}  & \textbf{83.18$^{e}$} & Out-of-distrib. flag & $\sim$80.1 \\
\textbf{Latency ($\mu$s)} & \textbf{142.2$^{a}$ / 161.9$^{b}$} & $\sim$12.5          & $\sim$283 \\
\textbf{Power (W)}      & \textbf{0.290$^{c}$ / 1.98$^{d}$} & $\sim$1.5 & $\sim$3.86 \\
\textbf{Energy/Inf.}    & \textbf{41.2~$\mu$J$^{c}$}    & $\sim$18.8~$\mu$J     & $\sim$1.09~mJ \\
\textbf{Weights Footprint} & \textbf{28.6~KB}     & $\sim$50~KB           & $\sim$187.5~KB \\
\textbf{Target Hardware}& \textbf{Zynq-7000} & Kintex-7 & {Zynq-7000} \\
\bottomrule
\end{tabular}

\vspace{2pt}
\parbox{\columnwidth}{\scriptsize $^{a}$compute, RTL simulation at 142.86~MHz; $^{b}$measured on ZedBoard, end to end incl.\ AXI input transfer; $^{c}$RadarNet-1D accelerator (PL), Vivado estimate; $^{d}$whole chip incl.\ PS7; $^{e}$integer model, executed bit-exactly by the hardware.}
\end{table}
%%%%%%%%%%%%%%%%%%%%%%%%%%%%%%%
%Total power is 0.671 W, composed of 0.559 W dynamic power (83\%) and 0.113 W device static power (17\%). The power consumed by clock distribution and routing is 0.271 W (49\%) and 0.152 W (27\%), respectively, which contributes to the highest power portion of the dynamic power. Combinational logic power consumes 0.130 W (23\%), and the power of BRAM and I/O is negligible.  The resulting junction temperature of 32.7$^{\circ}$C gives 52.3$^{\circ}$C of margin above the maximum junction temperature expected for systems deployed in the edge.
% [FIX] Added latency breakup: 56 cycles at 142 MHz design clock = 0.56 μs, at achieved Fmax 142 MHz = 0.776 μs
The implemented system closes timing at 142.86~MHz with a worst negative slack (WNS) of $+0.385$~ns and a worst hold slack of $+0.053$~ns. As shown in Table~\ref{tab:utilization}, the design uses 5.86\% of the total LUTs, 2.30\% of the total flip-flops, and 47.27\% of the DSP slices. This indicates that the proposed ResNet-1D inference datapath is a good fit for the target FPGA logic fabric. The BRAM utilization of 10.00\% indicates that both the weight ROM and intermediate feature maps can be comfortably stored in on-chip memory resources, which aligns with our 28.6~KB INT8 weight storage strategy. The lower half of Table~\ref{tab:power} breaks the latency down by component. The seven convolutions account for 97.9\% of the 20,315 compute cycles, of which 19,720 cycles (97.1\%) issue a MAC to the PE array; the last residual convolution alone takes 30\%, since it has the largest $C_{in}K$ product. On the board, the 128 I/Q samples arrive as individual AXI4-Lite writes from the ARM core, which adds 18.9~$\mu$s, so the measured end-to-end latency is 161.9~$\mu$s against 9.787~ms for the same integer model running on the 666.7~MHz Cortex-A9, a 60.5$\times$ speed-up with bit-exact logits on all 50 test vectors. %Input/output (IO) usage of 0.00\% is also minimal due to the fully self-contained accelerator running from internal BRAM memory.

Fig.~\ref{fig:waveform} illustrates the behavioral simulation waveform for the accelerator executing consecutively across test vectors from the 50-vector testbench. Each inference begins with a \texttt{start} strobe followed by 128 signed INT8 I/Q sample pairs on \texttt{iq\_data[15:0]}; \texttt{busy} then stays asserted while the descriptor sequencer time-multiplexes the 64-MAC/cycle PE array across all eight network layers. On inference completion, \texttt{done} and \texttt{result\_valid} pulse high, showing the predicted operational category on \texttt{class\_id[1:0]} and the complete set of four 32-bit classification logits on \texttt{logits[127:0]}.% With a 10~ns clock period in simulation, the measured duration between successive inference completions is 203.15~$\mu$s, directly corroborating the cycle count of 20,315 clock cycles ($20,315 \times 7.00\text{~ns} = 142.2\text{~}\mu\text{s}$ at 142~MHz system operating clock).

% Placeholder for Vivado Schematic
%\begin{figure}[htbp]
%\centerline{\includegraphics[width=\columnwidth]{schematic.png}}
%\caption{Vivado Schematic --- Top-level netlist showing VIO, radarnet\_top, ILA, LUT1, and dbg\_hub blocks (7 Cells, 1 I/O Port, 173 Nets).}
%\label{fig:schematic}
%\end{figure}
%%%%%%%%%%%%%%%%%%%%%%%%%%%%%%

%%%%%%%%%%%%%%%%%%%%%%%%%%%%%%
% Placeholder for Simulation Waveform
%\begin{figure}[htbp]
%\centerline{\includegraphics[width=\columnwidth]{behavorialsim.png}}
%\caption{Vivado Simulation Waveform --- Behavioral simulation showing incrementing cycle\_count every inference, FSM state transitions, I/Q data streaming, and control signals.}
%\label{fig:waveform}
%\end{figure}
%\subsection{Comparison with Prior FPGA RF Classifiers}

%The proposed accelerator is contextualized relative to other FPGA-based RF classification accelerators in Table~\ref{tab:comparison}.
%DeepRadio, an FPGA-based deep learning-based RF signal classifier, was proposed by Chang \emph{et al.}~\cite{b8} that can classify the signals received from the RF frontend in real-time while consuming very little power. The framework proposed by Umuroglu \emph{et al.}~\cite{b14}, FINN, leverages binary neural networks implemented on Xilinx FPGAs to deliver extremely high throughput while sacrificing the expressiveness of the model. hls4ml is another framework proposed by Duarte \emph{et al.}~\cite{b12} that automatically converts trained neural networks into FPGA firmware using HLS, and BatchNorm fusion support was available.
%%%%%%%%%%%%%%%%%%%%%%%%%%%%%%

%%%%%%%%%%%%%%%%%%%%%%%%%%%%%%
Table~\ref{tab:comparison} compares the proposed accelerator with FPGA-based RF classifiers and anomaly detectors that were previously implemented on comparable hardware platforms. Unlike the unsupervised autoencoder in~\cite{b10}, which is restricted to flagging out-of-distribution anomalies without identifying operational categories, our design performs explicit multi-class operational threat classification. When compared against the Zynq-7000 binarized neural network baseline implemented in~\cite{b14}, which suffers noticeable accuracy loss on continuous I/Q samples due to extreme 1-bit quantization, our integer-only INT8 datapath stays within 0.08 percentage points of floating-point accuracy while achieving 1.75$\times$ lower end-to-end latency (161.9~$\mu$s vs. $\sim$283~$\mu$s) and about 26$\times$ lower energy per inference (41.2~$\mu$J vs. $\sim$1.09~mJ).

\section{Conclusion}
This work presented an end-to-end system design for real-time RF anomaly detection at the sensor edge. By combining algebraic BN folding, per-channel INT8 quantization, and fixed-point scale-aligned residual additions, the 8-layer ResNet-1D is compressed to 28.6~KB of on-chip memory while maintaining 83.18\% test accuracy. Implementation on a Xilinx Zynq-7000 ZedBoard at 142.86~MHz shows that the accelerator achieves a measured end-to-end latency of 161.9~$\mu$s, 60.5$\times$ faster than the on-chip ARM core, with an estimated 0.290~W of PL power and 5.86\% LUT utilization, demonstrating its suitability for deployment in edge-based spectrum-monitoring systems.
%, which is better by over four orders of magnitude compared with GPU-based solutions.
%These characteristics make the proposed accelerator a strong candidate for deployment in edge-based spectrum-monitoring systems where low latency, low power, and deterministic operation are critical.

%\section{Acknowledgments}
%Portions of this manuscript were refined using AI-based writing assistants for grammar, clarity, and style. All scientific content, methodology, and analysis were developed and authored by the researchers.

\begin{thebibliography}{00}

\bibitem{b1} A. C. S. Thakur and M. Imtiaz, ``Deep learning for automatic modulation classification: A review,'' \emph{Electronics}, vol. 15, no. 10, p. 2163, May 2026.

\bibitem{b2} H. Xu, Y. Liang, R. Xie, and Y. Kong, ``The deep learning evolution in wireless physical layer communications: Applications, challenges, and evolutionary directions,'' \emph{Sensors}, vol. 26, no. 11, p. 3609, Jun. 2026.

\bibitem{b3} DeepSig Inc., ``RadioML 2018.01A dataset,'' 2018. [Online]. Available: https://www.deepsig.ai/datasets

\bibitem{b4} S. Ansari \emph{et al.}, ``Attention-enhanced hybrid automatic modulation classification for advanced wireless communication systems: A deep learning-transformer framework,'' \emph{IEEE Access}, vol. 13, 2025.

\bibitem{b5} F. Zhu \emph{et al.}, ``Robust deep learning-based physical layer communications: Strategies and approaches,'' \emph{IEEE Netw.}, 2025.

\bibitem{b6} S. Chatterjee, S. Ghosh, T. Ghosh, and H. Rahaman, ``Architectural design and performance analysis of FPGA based AI accelerators: A comprehensive review,'' arXiv preprint arXiv:2603.08740, Mar. 2026.

\bibitem{b7} X. Chen, X. Tian, Q. Zheng, and W. Wang, ``Attention mechanism enabled wireless modulation recognition network,'' \emph{Eng. Lett.}, vol. 33, no. 4, Apr. 2025.

\bibitem{b8} E. Chang, D. Cheng, and F. T. Luk, ``DeepRadio: Real-time FPGA deep learning for RF signal classification in electronic warfare,'' in \emph{Proc. IEEE Mil. Commun. Conf. (MILCOM)}, Norfolk, VA, Nov. 2019, pp. 1--6.

\bibitem{b9} E. Chang \emph{et al.}, ``Real-time modulation classification on Xilinx RFSoC using deep neural networks,'' in \emph{Proc. IEEE Int. Symp. Circuits Syst. (ISCAS)}, Seville, Spain, May 2020, pp. 1--5.

\bibitem{b10} D. J. M. Moss, D. Boland, P. Pourbeik, and P. H. W. Leong, ``Real-time FPGA-based anomaly detection for radio frequency signals,'' in \emph{Proc. IEEE Int. Symp. Circuits Syst. (ISCAS)}, Florence, Italy, May 2018, pp. 1--5.

\bibitem{b11} T. J. O'Shea, T. Roy, and T. C. Clancy, ``Over-the-air deep learning based radio signal classification,'' \emph{IEEE J. Sel. Topics Signal Process.}, vol. 12, no. 1, pp. 168--179, Feb. 2018.

%\bibitem{b12} J. Duarte \emph{et al.}, ``Fast inference of deep neural networks in FPGAs for particle physics,'' \emph{J. Instrum.}, vol. 13, no. 07, p. P07027, Jul. 2018.

%\bibitem{b13} X. Wang, Y. Zhao, and Z. Huang, ``A survey of deep transfer learning in automatic modulation classification,'' \emph{IEEE Trans. Cogn. Commun. Netw.}, vol. 11, no. 3, pp. 1357-1381, June 2025.

%\bibitem{b14} L. Zhang, T. Zhou, J. Yang, Y. Li, Z. Zhang, X. Hu, and Y. Peng, ``A real-time signal measurement system using FPGA-based deep learning accelerators and microwave photonic,'' \emph{Remote Sens.}, vol. 16, no. 23, p. 4358, Nov. 2024.

%\bibitem{b15} K. He, X. Zhang, S. Ren, and J. Sun, ``Deep residual learning for image recognition,'' in \emph{Proc. IEEE Conf. Comput. Vis. Pattern Recognit. (CVPR)}, Las Vegas, NV, Jun. 2016, pp. 770--778.

%\bibitem{b16} Y. LeCun, L. Bottou, Y. Bengio, and P. Haffner, ``Gradient-based learning applied to document recognition,'' \emph{Proc. IEEE}, vol. 86, no. 11, pp. 2278--2324, Nov. 1998.

%\bibitem{b17} S. Ioffe and C. Szegedy, ``Batch normalization: Accelerating deep network training by reducing internal covariate shift,'' in \emph{Proc. 32nd Int. Conf. Mach. Learn. (ICML)}, Lille, France, Jul. 2015, pp. 448--456.

\bibitem{b12} B. Jacob \emph{et al.}, ``Quantization and training of neural networks for efficient integer-arithmetic-only inference,'' in \emph{Proc. IEEE CVPR}, Salt Lake City, UT, Jun. 2018, pp. 2704--2713.

\bibitem{b13} D. P. Kingma and J. Ba, ``Adam: A method for stochastic optimization,'' in \emph{Proc. Int. Conf. Learn. Representations (ICLR)}, San Diego, CA, USA, May 2015.

\bibitem{b14} Y. Umuroglu \emph{et al.}, ``FINN: A framework for fast, scalable binarized neural network inference,'' in \emph{Proc. ACM/SIGDA Int. Symp. Field-Program. Gate Arrays (FPGA)}, Monterey, CA, Feb. 2017, pp. 65--74.

\end{thebibliography}

\end{document}
